% status: ZERO
% Chapter 37 of the second edition. Plan: docs/STRUCTURE.md. Rules: AUTHORING.md.
% Measurements: research/measurements/2026-09-24-m1-part7.md (contention, gpucopy,
% mmapload, offload, repack, metalbench); Hermon's own records at commit 2a3fd52.
\chapter{Unified Inference Memory}\label{ch:unified-memory}

\begin{ChapterIntent}
When the CPU and GPU share one physical memory, the tiers of the previous
chapters collapse into one pool with one bandwidth. This chapter examines
unified memory on Apple silicon and on coherent CPU--GPU superchips, what it
gives up relative to dedicated HBM, and the design of a single allocator and
residency policy that manages weights, KV blocks and experts together rather
than as three competing caches.
\end{ChapterIntent}

\OpeningSection{One memory, two processors}

On the M1, llama.cpp decoded Llama~3.2~3B at 24.8 to 25.6 tokens per second on
the GPU. Then one ordinary CPU process began copying a 256~MiB array in a loop,
and decoding fell to 13.5 tokens per second. With two such processes it ran at
13.0; with four, at 14.2 (measured, interleaved with runs that had no copying
processes; record \texttt{2026-09-24-m1-part7}). Nothing about the GPU's work
had changed, and the copying processes used none of the GPU. They used the
memory. A GPU copy measured the same way moved 49 to 56~GB/s alone; with one CPU
process copying beside it, the GPU got 35.6~GB/s and the CPU 29.7 --- together
65.3~GB/s, more than either drew alone. The limit was the memory, not either
processor.

\EngineFigure{ch37-shared-bandwidth}{Memory bandwidth shared between the M1's GPU
and CPU. Left: a bulk copy on the GPU, alone and beside one, two or four CPU
processes copying memory (measured). Right: llama.cpp decoding Llama~3.2~3B
beside the same processes; decode's reads are its token rate times the 2.02~GB
model (derived), the CPU processes' rates are measured. The GPU keeps about
two thirds of its bulk-copy rate and decode about half its token rate however
many processes compete: one CPU process is enough to take the rest. The
zero-process bars average four interleaved baseline runs.}{fig:ch37-shared-bandwidth}
On a machine with one memory, the tiers of \Chref{offloading} collapse. There
is no link to cross and no copy to make: the GPU reads the model file's pages
where the operating system put them. What remains is one pool of capacity and
one pool of bandwidth, shared by every processor and every program on the
machine --- and a design question the previous chapters could avoid: who
decides what stays in it.

\Foundations{Shared addresses do not imply free access}{
An address space describes what a processor can name; physical placement and the
memory fabric determine what it costs to access. CPU and GPU may share physical
memory, or a runtime may migrate pages behind a unified-address interface. Those are
different mechanisms. Even physically shared memory has finite bandwidth and competing
users. Synchronization still determines when a producer's writes are usable by a
consumer. Chapters~\ref{app:gpu} and~\ref{app:cpu} describe the two execution
agents. Chapter~\ref{ch:offloading} explains the transfers that a particular
shared-memory design does, or does not, eliminate.
}

\NapkinSection{Napkin math: one bandwidth shared three ways}

Decode reads its weights once per step, so on a machine whose memory delivers
$\beta$ in total, of which other processors take $\beta_{\mathrm{other}}$, a
model that reads $W$ bytes per token decodes at most
\begin{equation}
r \le \frac{\beta - \beta_{\mathrm{other}}}{W}
\label{eq:shared-bw}
\end{equation}
tokens per second. Llama~3.2~3B in Q4\_K\_M form is a 2.02~GB file, and because
its output projection shares the embedding matrix a decode step reads nearly
all of it. At 25.5 tokens per second the GPU drew about 51~GB/s, close to the
49 to 56~GB/s of its bulk copy alone. With a CPU process copying beside it the
bulk copy fell to 35.6~GB/s, which by Equation~\ref{eq:shared-bw} bounds decode
at 17.6 tokens per second; decode measured 13.5 (derived from the measured
rates). The bound held, and decode lost more than the bulk copy did --- 47\%
against 35\% --- plausibly because a decode step is a chain of small dependent
kernels, each exposed to the contention, where a bulk copy keeps many requests
in flight.

The same equation prices the machines built around unified memory. Apple does
not publish the M1's bandwidth; it describes the M5's 153~GB/s as more than
twice the M1's \cite{apple2025m5}. A 4-bit model of 70 billion parameters reads
about 40~GB per token, so its decode ceiling is 3.8 tokens per second at
153~GB/s, 6.8 on a 273~GB/s DGX~Spark \cite{nvidia2026dgxspark}, 15 on an M5~Max
at 614~GB/s and 30 on an M5~Ultra at 1.2~TB/s \cite{apple2026m5ultra} --- before
the cache, and before anything else on the machine takes its share
(derived). The same model on a data-center GPU's 3.35~TB/s of HBM would reach
84 \cite{nvidia2026h100}; unified memory buys capacity and pays in bandwidth.

\section{One pool instead of tiers}

\Chref{offloading} priced every byte by the tier it lived on and the link it
crossed. With one memory, a byte's price depends only on which processor reads
it and how fast that processor can read. On the M1, moving all of a Qwen3-8B-architecture
model's layers from the GPU to the CPU moved no bytes: decoding went from 10.0
to 10.7 tokens per second to 5.9 to 8.0, because the four CPU cores read the
same memory at about two thirds of the GPU's rate (measured). The capacity
is pooled too. The operating system, the user's applications, the model's
weights, its cache and any expert cache all draw on the same 16~GB, and
nothing but the pager arbitrates between them. \Chref{kv-storage-tier} found
what that means for a cache that can be paged out; this chapter's measurements
found it again for weights.

Pooling also means that a second copy costs the same scarce memory as the
first. When llama.cpp splits a model between the processors it keeps a
repacked copy of the CPU's layers while the GPU maps the whole file
(\Chref{offloading}). With 9 of 37 layers on the GPU that copy was 2.9~GiB,
and on this machine it meant 6.0 to 22.7~GiB of swap-outs during one short run
and decoding at 0.4 to 1.4 tokens per second, against 5.7 to 9.1 with the copy
turned off (measured). Without repacking, a model run entirely on the CPU
existed only as pages of the mapped file, as resident as the operating system
chose to keep them. In one run on the overcommitted machine it evicted them:
decoding ran at 0.10 tokens per second while 161~GiB were paged in from the SSD
in five and a half minutes --- the model read from disk token by token. Two days
later, with less pressure on memory, the same configuration kept its pages and
decoded at 8.0 (measured). Capacity on one memory is not a
number to fill; it is a budget the engine has to defend.

\section{Apple silicon}

Metal gives a program two ways to put data in front of the GPU. It can
allocate a buffer and copy data into it, or it can wrap memory the program
already has --- page-aligned, in shared storage mode --- with
\texttt{newBufferWithBytesNoCopy}, so that the buffer's contents are the
program's own pages. llama.cpp takes the second path for model weights: it maps
the GGUF file into memory and wraps the mapping, splitting it into overlapping
views when it exceeds the device's largest buffer
(\texttt{ggml/src/ggml-metal/ggml-metal-device.m} at \texttt{d006858}; DEFAULT).
The server's log shows the result: for the Qwen3-8B-architecture model a single
4,977.62~MiB \texttt{MTL0\_Mapped} buffer covers all of the file's tensor data,
while with \texttt{--no-mmap} it allocates a 4,643.78~MiB buffer and reads the
weights into it (measured). Once the file was in the page cache, both started in
4.6~seconds; the copy is not what a laptop waits for at startup.

Wrapping is free; keeping the pages resident is not. On macOS~15 and later,
llama.cpp adds every buffer to a Metal residency set and, from a background
thread, requests residency again every half second until three minutes after
the last use, which keeps the weights resident through short pauses between
requests (same file; DEFAULT, disabled by
\texttt{GGML\_METAL\_NO\_RESIDENCY}). The device also reports how much it
recommends a program keep in use at once: on the 16~GB M1,
\texttt{recommendedMaxWorkingSetSize} is 12,713~MB (measured, from the log).
Newer chips raise both sides of Equation~\ref{eq:shared-bw}: the M5 Max
supports 128~GB at up to 614~GB/s, and the M5 Ultra 512~GB at 1.2~TB/s, which
Apple markets for running models with hundreds of billions of parameters on
the device \cite{apple2026m5ultra}.

\section{Coherent CPU--GPU superchips}

Data-center designs reach a similar place from the other direction. NVIDIA's
Grace Hopper keeps two memories --- LPDDR5X on the CPU and HBM on the GPU ---
and joins them with NVLink-C2C at 900~GB/s, 450 in each direction, which NVIDIA
compares to seven times PCIe Gen~5. The link is cache-coherent: CPU and GPU
threads access each other's memory without page migration, through the CPU's
address translation \cite{nvidia2022gracehopper,nvidia2026gh200}. The tiers of
\Chref{offloading} are still there, but the link between them is fast enough
that host memory becomes a place to keep cache and weights that are read less
often, rather than a last resort. DGX~Spark's GB10 is closer to Apple's model:
128~GB of coherent unified LPDDR5x memory at 273~GB/s, shared by a 20-core Arm
CPU and a Blackwell GPU \cite{nvidia2026dgxspark}. AMD's Instinct MI300A goes
furthest: 24 CPU cores and 228 GPU compute units share one 128~GB pool of HBM3,
fully coherent, so the unified pool runs at HBM speed \cite{amd2026mi300a}.

These designs put different points on one curve. Unified LPDDR memory gives
large capacity at modest bandwidth; HBM gives the most bandwidth and the least
capacity; a coherent link lets a system hold both and choose per byte. Which one a workload wants is Equation~\ref{eq:shared-bw} again, with
capacity as the constraint that decides whether the model runs at all.

\section{One allocator for weights, cache and experts}

Every engine in this book manages its memory as separate budgets. Weights are
loaded once and stay; the KV cache gets a fixed allocation, carved into blocks
(\Chref{paged-kv}); an expert cache, where there is one, gets a fixed size of
its own (\Chref{expert-streaming}); and the operating system's page cache sits
under all three with a policy of its own. On one memory those budgets compete
for the same bytes, and a static split is wrong for every workload but the one
it was tuned on: a long-context request wants cache, an expert-heavy one wants
experts, and neither can borrow from the other.

Hermon's kernel library puts KV blocks and expert records in the same block
pool, with the same reference counts and free list --- one arena, with the
axis changing from tokens to experts --- so that the two could in principle be
traded (LIBRARY at \texttt{2a3fd52}). The policy that would trade them does not
exist. Hermon's own storage document is plain about it: its three tiers ---
device memory, host memory, NVMe --- run three unrelated policies with no
promotion or demotion between them; a KV block under pressure cannot spill to
the drive, it fails; and the design's target of a 20\% gain from unified KV and
expert budgeting on a mixed workload is unrun (DESIGN;
\texttt{docs/STORAGE\_ARCHITECTURE.md} and \texttt{docs/KERNEL\_DESIGN.md} at
\texttt{2a3fd52}).

\section{A single residency policy}

What a unified policy must decide is which bytes stay resident, and the
previous chapters supply the value of each kind: the rate at which it is read
per byte of memory it occupies. Dense weights are read by every step. The
active cache of a running sequence is read by every step of that sequence. An
expert is read with the probability its router gives it. The cache of a paused
conversation is read by no step until its user returns, and a prefix by the
requests that share it. Ranking bytes by reads per second per byte, and
charging each tier's miss cost from Equation~\ref{eq:offload-time}, gives one
ordering for all of them --- the same economics as a buffer pool in a database,
with one difference that matters: an inference step's reads are known a layer
ahead for weights and cache, so prefetching them can be exact, and experts can
be predicted a layer ahead, imperfectly (\Chref{expert-streaming}).

Two measured failures bound the design. WARP found that an expert cache larger
than physical memory ran eight times slower, because the operating system paged
it out and every hit became a page fault \cite{sqliteai2026warp}. And on the
M1, llama.cpp's host-memory prompt cache took 67 seconds to return a cache that
a file restored in 4 (\Chref{kv-storage-tier}). An engine on unified memory
must hold its budget below what the operating system will leave it, or its
residency policy is only a suggestion.

\section{Contention for shared bandwidth}

The opening's measurement is Equation~\ref{eq:shared-bw} in action. On one
memory, bandwidth is the resource every processor shares, and nothing
partitions it: the CPU copy that halved decoding was an ordinary process with
no priority. On the M1 one copying process alone moved 55~GB/s, and two or four
moved no more in total --- 57 and 53~GB/s --- because a single process can
saturate the memory; with the GPU also copying, the pair reached 55 to 65~GB/s
between them (measured). For an engine this has two consequences. A laptop's
token rate depends on what else the laptop is doing --- this book's own
measurements on the M1 moved by factors of two and more with the machine's
load --- and
work an engine moves to the CPU to relieve the GPU competes with the GPU for
the one resource decode needs most.

A common address does not imply that two processors may read and write without coordination. The useful change in this design is that weights and state can reside in one physical pool. The allocator still tracks last use, and concurrent readers still compete for memory bandwidth.

\EngineFigure{sys-shared-memory}{Component layout of a physically shared-memory system, not merely a unified virtual address space. CPU and GPU retain distinct execution and caching behavior while sharing a capacity and traffic budget. Synchronization is still required at producer/consumer boundaries.}{fig:sys-shared-memory}

\LedgerSection

Running inference on one memory shared by CPU and GPU:

\begin{ConstraintLedger}
Memory capacity & buys & One pool as large as the machine's memory, with no copies between CPU and GPU. \\
Memory bandwidth & spends & One bandwidth, usually far below HBM's, shared by every processor and program; decode slows when anything else reads memory. \\
Latency & risks & The operating system can page out weights or cache that the engine considers resident. \\
Compute & moves & Work moved to the CPU competes with the GPU for bandwidth, not only for arithmetic. \\
\end{ConstraintLedger}

\LabSection{one allocator, three kinds of data}

\LabIntent{In \texttt{code/labs/ch37-unified-memory/}, one block pool that
holds weight pages, KV blocks and expert records under a single budget, with a
residency policy that ranks blocks by reads per second per byte. It replays a
mixed trace --- long-context requests, paused conversations and MoE routing ---
and compares the unified policy with the same budget split statically three
ways, reporting bytes read from the slower tier and a predicted token rate
from Equation~\ref{eq:shared-bw}. Every read is checked against a content hash,
so a policy bug cannot hide as a speed-up. Break it by giving the pool more
than the machine will leave it and measure the page faults.}

\HappenedSection{zero-copy on unified memory}

Hermon's Metal attention kernel was built on one measured fact: on the M1,
\texttt{newBufferWithBytesNoCopy} wraps the KV arena's own pages, so a GPU
kernel writing through the buffer changes the CPU's memory in place, and the
cache is never staged. Without that, the developers computed, a decode step
would copy the whole context every token --- 512~MiB for Llama~3 8B at 4,096
positions. The library verifies the alias when it attaches a pool rather than
assuming it (Hermon \texttt{docs/KERNEL\_DESIGN.md} section K4; the developers'
measurement).

Zero-copy removed the copy and did not make the GPU win. Against the CPU's
vector kernel, the Metal kernel lost below about 4,000 tokens of context and
won by less than two times above it, because each call paid a fixed cost of
about 0.94~ms: a command-buffer round trip of 0.19~ms that could not be
removed, plus launch and a grid too small to fill the GPU at decode. The layers
could not share one command buffer, because the host computed everything
between two attention calls, so a forward pass paid the round trip once per
layer. An attempt to save bandwidth by fusing the query heads that share a
key/value head was bit-identical and slower in every shape measured: the reads
it saved had already been served from the GPU's cache, and it cost the one
thing the kernel lacked, parallelism (section K4.1). The kernel shipped behind a
gate that is a pure function of the problem's shape, so the same request always
takes the same path. Re-run at the same commit on this book's M1, the benchmark
reproduced the developers' crossover within a few percent: with 32 query and 8
key/value heads, Metal took 0.21, 0.44 and 0.79 times the CPU thread's speed at
128, 512 and 2,048 tokens of context and won by 1.56 and 1.79 times at 8,192 and
32,768, against the developers' 0.19 to 1.81 (measured; record
\texttt{2026-09-24-m1-part7}). Unified memory makes sharing free; it does not
make crossing between processors free.

\FrontierSection{2026-09-24}

Unified memory has become a product category for local inference: Apple's
desktops now reach 512~GB at 1.2~TB/s \cite{apple2026m5ultra}, NVIDIA sells a
128~GB unified desktop system \cite{nvidia2026dgxspark}, and data-center parts
join CPU and GPU memory coherently
\cite{nvidia2022gracehopper,amd2026mi300a}. Engines have been slower to follow.
Their allocators still divide memory into fixed budgets for weights, cache and
experts, and leave the rest to the operating system. The open problem is the
one this chapter's lab poses: a single residency policy that ranks weights,
cache blocks and experts by their value per byte, owns its budget against the
operating system, and prefetches using what the model's structure reveals
about the next layer.

\ExercisesSection

\begin{enumerate}
\item \emph{Derivation.} Extend Equation~\ref{eq:shared-bw} to a batch of $B$
  sequences with $S$ tokens of cache each, and find the batch at which the
  cache reads equal the weight reads for Qwen3-8B on a 273~GB/s machine.
\item \emph{Prediction.} A 400-billion-parameter MoE model with 17 billion
  active parameters is stored at 4 bits on a 512~GB, 1.2~TB/s machine. Predict
  its decode ceiling at batch one, and the batch at which its reads of expert
  weights stop growing (\Chref{mixture-of-experts}).
\item \emph{Implementation.} Add a memory-copying background process to the
  lab's measurement and verify Equation~\ref{eq:shared-bw} on your machine at
  three levels of contention.
\item \emph{Falsification.} Find a workload for which a static three-way split
  of the pool beats the lab's unified policy. What does the unified policy fail
  to know?
\end{enumerate}
